\section{Реализация веб-приложения}

\subsection{Выбор архитектуры и структуры проекта}

Практическая реализация веб-приложения выполнена в соответствии с проектными
решениями, приведёнными во второй части курсового проекта. Система построена по
клиент-серверной архитектуре: пользователь взаимодействует с интерфейсом
React-приложения, клиентская часть отправляет HTTP-запросы к серверу Django REST
Framework, а серверная часть выполняет бизнес-логику и сохраняет данные в
PostgreSQL.

Выбранная архитектура позволяет разделить ответственность между частями системы:
клиентская часть отвечает за отображение интерфейса, обработку форм и маршрутизацию
между страницами, серверная часть отвечает за проверку данных, авторизацию,
работу с бронированиями, формирование PDF-подтверждений, отправку email-уведомлений
и экспорт отчётов. Для воспроизводимого запуска проекта используются Docker и
файл \texttt{docker-compose.yml}.

Структура реализованного проекта представлена в листинге~\ref{lst:project_tree}.

\begin{lstlisting}[caption={Структура реализованного проекта}, label={lst:project_tree}]
reservation_service/
  backend/
    manage.py
    requirements.txt
    reservation_service/
      settings.py
      urls.py
    booking/
      models.py
      serializers.py
      views.py
      services.py
      urls.py
      admin.py
      tests.py
  frontend/
    package.json
    src/
      App.jsx
      api/client.js
      components/Layout.jsx
      pages/
      styles/style.css
  docker-compose.yml
  .env.example
  paper.tex
  ch1.tex
  ch2.tex
  ch3.tex
\end{lstlisting}

\subsection{Реализация серверной части}

Серверная часть разработана на базе фреймворка Django и библиотеки Django REST
Framework. В приложении \texttt{booking} размещены модели данных, сериализаторы,
представления API, сервисные функции и настройки административной панели.

Основными сущностями системы являются: пользователь, столик, позиция меню,
бронирование, уведомление и файл бронирования. Для пользователя используется
стандартная модель \texttt{User}, входящая в Django. Разграничение доступа между
обычным авторизованным пользователем и администратором выполняется через поле
\texttt{is\_staff}. Обычный пользователь может создавать и просматривать только
свои бронирования, администратор получает доступ к управлению столиками, меню,
бронированиями и отчётами.

Пример модели бронирования представлен в листинге~\ref{lst:booking_model}.

\begin{lstlisting}[caption={Фрагмент модели бронирования}, label={lst:booking_model}]
class Booking(models.Model):
    user = models.ForeignKey(User, on_delete=models.CASCADE)
    table = models.ForeignKey(RestaurantTable, on_delete=models.PROTECT)
    date = models.DateField()
    time_start = models.TimeField()
    time_end = models.TimeField()
    guests_count = models.PositiveIntegerField(default=1)
    status = models.CharField(max_length=20, default='pending')
    pdf_path = models.CharField(max_length=255, blank=True)
\end{lstlisting}

Для передачи данных между клиентом и сервером используются сериализаторы. Они
преобразуют объекты Django ORM в JSON и выполняют дополнительную проверку данных.
Например, при создании бронирования проверяется корректность временного интервала,
соответствие количества гостей вместимости столика и отсутствие пересечений с уже
существующими бронированиями.

\subsection{Реализация базы данных}

В качестве основной базы данных в проекте предусмотрена PostgreSQL. При запуске проекта через
Docker создаётся отдельный сервис \texttt{db}, к которому подключается backend.
В процессе разработки также предусмотрен резервный вариант с SQLite, если переменные
окружения PostgreSQL не заданы. Для развёртывания демонстрационной версии приложения
на бесплатном тарифе PythonAnywhere использовалась SQLite, так как она не требует
отдельного серверного процесса и подходит для учебной демонстрации функциональности.

В таблице~\ref{tab:db_entities} приведены основные таблицы базы данных и их
назначение.

\begin{longtable}{|L{6.2cm}|L{7.8cm}|}
\caption{Основные таблицы базы данных}\label{tab:db_entities}\\
\hline
\textbf{Таблица} & \textbf{Назначение} \\
\hline
\endfirsthead
\hline
\textbf{Таблица} & \textbf{Назначение} \\
\hline
\endhead
\texttt{auth\_user} & хранение учётных записей пользователей и администраторов \\
\hline
\texttt{booking\_restauranttable} & хранение столиков, вместимости и изображений \\
\hline
\texttt{booking\_menuitem} & хранение позиций меню, категорий, цен и изображений блюд \\
\hline
\texttt{booking\_booking} & хранение бронирований, даты, времени, статуса и PDF-подтверждения \\
\hline
\texttt{booking\_notification} & хранение внутренних и email-уведомлений пользователя \\
\hline
\texttt{booking\_bookingfile} & хранение PDF-файлов, связанных с конкретным бронированием \\
\hline
\end{longtable}

Связь таблиц реализована через внешние ключи. Бронирование связано с пользователем
и столиком, уведомление связано с пользователем и, при необходимости, с конкретным
бронированием, а файл подтверждения связан с записью бронирования.

\subsection{Реализация аутентификации пользователей}

Для входа в систему используется встроенная система пользователей Django и токенная
аутентификация Django REST Framework. После регистрации или входа сервер возвращает
клиентской части токен. React-приложение сохраняет токен в локальном хранилище и
передаёт его в заголовке \texttt{Authorization} при последующих запросах к API.

Фрагмент клиентского кода, отвечающий за добавление токена к запросам, показан в
листинге~\ref{lst:token_client}.

\begin{lstlisting}[caption={Добавление токена авторизации к запросам}, label={lst:token_client}]
api.interceptors.request.use((config) => {
  const token = localStorage.getItem('token');
  if (token) {
    config.headers.Authorization = `Token ${token}`;
  }
  return config;
});
\end{lstlisting}

На стороне frontend реализованы страницы регистрации и входа. После успешного входа
пользователь переходит к просмотру меню и форме бронирования. Если пользователь
является администратором, в интерфейсе отображается ссылка на административные
разделы.

\subsection{Реализация бронирования и проверки доступности столиков}

Одной из ключевых функций системы является проверка доступности столика на выбранные
дату и время. Проверка выполняется на сервере перед созданием бронирования. Система
ищет активные бронирования того же столика на ту же дату и проверяет пересечение
временных интервалов. Если пересечение найдено, создание бронирования блокируется.

Фрагмент функции проверки доступности приведён в листинге~\ref{lst:availability}.

\begin{lstlisting}[caption={Проверка доступности столика}, label={lst:availability}]
def is_table_available(table, date, time_start, time_end):
    query = Booking.objects.filter(
        table=table,
        date=date,
        status__in=['pending', 'confirmed'],
    ).filter(time_start__lt=time_end, time_end__gt=time_start)
    return not query.exists()
\end{lstlisting}

На клиентской стороне реализована форма бронирования. Пользователь указывает дату,
время начала, время окончания, количество гостей и выбирает столик. Перед отправкой
формы можно выполнить отдельную проверку доступности, после которой в интерфейсе
отображается сообщение о том, доступен ли столик для выбранного времени.

Для повышения наглядности выбора столика в форме бронирования предусмотрен
визуальный предпросмотр. После выбора номера столика справа от формы отображается
схема выбранного столика. Если администратор загрузил собственное изображение
столика, используется оно; если изображение не загружено, отображается стандартная
схема по номеру столика. Такой подход не требует изменения структуры базы данных,
так как изображение уже связано с сущностью столика через поле \texttt{image}, но
делает интерфейс более понятным для пользователя.

\subsection{Реализация административной панели}

Административная часть реализована на двух уровнях. Первый уровень - стандартная
административная панель Django, позволяющая работать с моделями напрямую. Второй
уровень - пользовательский административный интерфейс на React, который использует
REST API и отображает основные функции, предусмотренные проектом.

В административной части реализованы следующие возможности:

\begin{itemize}
  \item просмотр списка бронирований;
  \item подтверждение или отмена бронирования;
  \item управление столиками;
  \item управление позициями меню;
  \item загрузка изображений столиков и блюд;
  \item формирование отчётов о занятости столиков;
  \item экспорт отчётов в CSV и PDF.
\end{itemize}

Для защиты административных методов используется проверка роли пользователя.
Операции изменения данных доступны только пользователям с признаком \texttt{is\_staff}.

\subsection{Реализация уведомлений, PDF-подтверждений и email-рассылки}

После создания бронирования система формирует внутреннее уведомление, создаёт
PDF-подтверждение и отправляет email на адрес пользователя. Уведомление сохраняется
в базе данных и отображается в интерфейсе через страницу уведомлений и значок
колокольчика в верхней панели приложения. Email-уведомление содержит основные
параметры бронирования: дату, время, номер столика, количество гостей и текущий
статус.

Для хранения состояния рассылки в модели уведомления используются поля
\texttt{delivery\_channel}, \texttt{email\_status} и \texttt{sent\_at}. Это позволяет
отслеживать, было ли уведомление отправлено только внутри сайта или дополнительно
продублировано на электронную почту.

Фрагмент создания уведомления показан в листинге~\ref{lst:notification}.

\begin{lstlisting}[caption={Создание уведомления о бронировании}, label={lst:notification}]
notification = Notification.objects.create(
    user=booking.user,
    booking=booking,
    title='Booking created',
    message='Booking details were generated by the system.',
    delivery_channel='site_email',
)
\end{lstlisting}

В режиме разработки используется консольный backend отправки почты: письма выводятся
в лог сервера. Для реальной отправки достаточно указать SMTP-настройки в файле
\texttt{.env}. Такой подход позволяет не хранить пароль от почтового ящика в коде.

\subsection{Реализация клиентской части}

Клиентская часть реализована на React. Основной компонент \texttt{App.jsx} отвечает
за переключение страниц, хранение текущего состояния интерфейса и загрузку количества
непрочитанных уведомлений. Запросы к серверу вынесены в отдельный модуль
\texttt{api/client.js}. Это упрощает поддержку проекта, потому что базовый адрес API
и добавление токена авторизации настраиваются в одном месте.

В приложении реализованы страницы:

\begin{itemize}
  \item вход и регистрация;
  \item просмотр меню;
  \item создание бронирования;
  \item просмотр уведомлений;
  \item управление бронированиями;
  \item управление столиками;
  \item управление меню;
  \item отчёты и экспорт данных.
\end{itemize}

Интерфейс оформлен с помощью CSS и адаптирован для просмотра на экранах разной
ширины. В верхней панели отображается название приложения, навигация и значок
колокольчика с количеством непрочитанных уведомлений.

\subsection{Запуск и тестирование приложения}

Для запуска проекта используется Docker Compose. Перед запуском необходимо создать
файл \texttt{.env} на основе \texttt{.env.example}. После этого приложение запускается
командой, приведённой в листинге~\ref{lst:docker_run}.

\begin{lstlisting}[caption={Запуск приложения}, label={lst:docker_run}]
docker compose up --build
\end{lstlisting}

После запуска доступны следующие адреса:

\begin{itemize}
  \item frontend: \texttt{http://localhost:5173};
  \item backend api: \texttt{http://localhost:8000/api/};
  \item административная панель Django: \texttt{http://localhost:8000/admin/}.
\end{itemize}

Для проверки серверной логики были добавлены автоматические тесты Django REST
Framework. Тесты размещены в файле \texttt{booking/tests.py} и проверяют наиболее
важные функции приложения: регистрацию и вход пользователя, проверку доступности
столика, создание бронирования, защиту от пересечения бронирований, создание
уведомлений и PDF-подтверждений, права администратора, загрузку изображений столиков
и блюд меню, а также экспорт отчётов.

Запуск автоматических тестов выполняется командой, приведённой в
листинге~\ref{lst:run_tests}.

\begin{lstlisting}[caption={Запуск автоматических тестов backend}, label={lst:run_tests}]
docker compose exec backend python manage.py test
\end{lstlisting}

Основные тестовые сценарии приведены в таблице~\ref{tab:testing_scenarios}.

\begin{longtable}{|L{5.0cm}|L{5.2cm}|L{3.4cm}|}
\caption{Сценарии тестирования приложения}\label{tab:testing_scenarios}\\
\hline
\textbf{Проверяемая функция} & \textbf{Сценарий} & \textbf{Ожидаемый результат} \\
\hline
\endfirsthead
\hline
\textbf{Проверяемая функция} & \textbf{Сценарий} & \textbf{Ожидаемый результат} \\
\hline
\endhead
Аутентификация & Регистрация нового пользователя и последующий вход в систему & Сервер возвращает токен авторизации \\
\hline
Проверка доступности & Проверка столика при отсутствии пересечений по времени & Столик доступен \\
\hline
Защита от двойного бронирования & Создание бронирования на уже занятый временной интервал & Сервер возвращает ошибку валидации \\
\hline
Создание бронирования & Авторизованный пользователь создаёт бронирование & Создаются бронирование, уведомление и PDF-файл \\
\hline
Права доступа & Обычный пользователь пытается создать столик & Сервер запрещает операцию \\
\hline
Административное управление & Администратор создаёт столик и загружает изображение & Столик сохраняется с изображением \\
\hline
Управление меню & Администратор создаёт позицию меню и загружает фото блюда & Блюдо сохраняется с изображением \\
\hline
Отчёты & Администратор запрашивает CSV и PDF отчёты & Сервер возвращает файлы отчётов \\
\hline
\end{longtable}

Дополнительно было выполнено ручное тестирование клиентской части. Вручную проверялись
отображение страниц, работа формы бронирования, визуальный предпросмотр выбранного
столика, отображение уведомлений через значок колокольчика, административные разделы
и корректность загрузки изображений через интерфейс администратора. Такой подход
позволяет проверить не только серверную логику, но и соответствие интерфейса
проектным шаблонам.

В результате тестирования подтверждено, что реализованная система выполняет основные
функции, предусмотренные заданием курсового проекта.

\subsection{Выводы по разделу}
Исходный код разработанного веб-приложения размещён в репозитории GitHub по адресу \url{https://github.com/anastasiaantipovamax/reservation_service}. Репозиторий содержит серверную часть на Django REST Framework, клиентскую часть на React, настройки Docker Compose, а также файлы пояснительной записки. Веб-приложение развёрнуто на платформе PythonAnywhere и доступно по адресу \url{https://anastasiaantipova.pythonanywhere.com/}.

В ходе реализации была создана рабочая версия веб-приложения для автоматизации
бронирования столиков в ресторане. Серверная часть на Django REST Framework
обеспечивает хранение данных, аутентификацию, проверку доступности столиков,
формирование PDF-подтверждений, уведомления и экспорт отчётов. Клиентская часть на
React предоставляет интерфейс для авторизованного пользователя и администратора.
Основной сценарий запуска с PostgreSQL и Docker делает приложение воспроизводимым и
приближает его к реальному сценарию эксплуатации, а демонстрационное размещение на
PythonAnywhere с SQLite позволило показать работу проекта на бесплатном хостинге без
изменения пользовательской функциональности.
